\begin{abstract}
Automated scientific discovery promises to accelerate progress across scientific domains. However, developing and evaluating an AI agent's
capacity for end-to-end scientific reasoning is challenging as running real-world experiments is often prohibitively expensive or infeasible. In this work we introduce \env,
the first virtual environment for developing and benchmarking an agent's ability to perform complete cycles of novel scientific discovery.
\env contains a variety of different challenges, covering topics as diverse as radioisotope dating, rocket science, and proteomics, to encourage development of {\it general} discovery skills rather than task-specific solutions.
\env itself is an inexpensive, simulated, text-based environment (with optional 2D visual overlay).
  It includes 120 different challenge tasks, spanning eight topics each with three levels of difficulty and
several parametric variations.
Each task requires an agent to form hypotheses, design and run experiments, analyze results, and act on conclusions.
\env further provides three automatic metrics for evaluating performance, based on (a) task completion, (b) task-relevant actions taken, and (c) the discovered explanatory knowledge.
We find that strong baseline agents, that perform well in prior published environments, struggle on most \env tasks,
suggesting that \env captures some of the novel challenges of discovery, and thus that \env may help
accelerate near-term development and assessment of scientific discovery competency in agents.
Code available at \href{https://github.com/allenai/discoveryworld}{github.com/allenai/discoveryworld}.\footnote{Released under Apache-2.0 license.}
\end{abstract}
